% !TEX program = xelatex
\documentclass[11pt,a4paper]{article}
\usepackage[margin=2.2cm,top=2cm]{geometry}
\usepackage{fontspec}
\usepackage{polyglossia}
\setdefaultlanguage{turkish}
\usepackage{amsmath,amssymb}
\usepackage{graphicx}
\usepackage{booktabs}
\usepackage{enumitem}
\usepackage{multicol}
\usepackage[most]{tcolorbox}
\usepackage{fancyhdr}
\usepackage{titlesec}
\usepackage{xcolor}
\usepackage{needspace}
\usepackage[hidelinks]{hyperref}

\definecolor{anamavi}{RGB}{0,45,98}
\definecolor{cozumyesil}{RGB}{0,110,50}
\definecolor{uyarikirmizi}{RGB}{180,20,20}

\titleformat{\section}{\large\bfseries\color{anamavi}}{\thesection.}{0.5em}{}
\titleformat{\subsection}{\normalsize\bfseries\color{anamavi}}{\thesubsection}{0.5em}{}

\pagestyle{fancy}\fancyhf{}
\lhead{\small MMAT105 Genel Matematik}
\rhead{\small Hafta 1 -- Ders Notu}
\cfoot{\small\thepage}
\renewcommand{\headrulewidth}{0.4pt}

\newtcolorbox{kural}[1][Kural]{colback=anamavi!5,colframe=anamavi,title=#1,fonttitle=\bfseries,breakable,left=4pt,right=4pt,top=2pt,bottom=2pt}
\newtcolorbox{dikkat}[1][Dikkat]{colback=uyarikirmizi!5,colframe=uyarikirmizi,title=#1,fonttitle=\bfseries,breakable,left=4pt,right=4pt,top=2pt,bottom=2pt}
\newcounter{ornek}
\newenvironment{ornek}{\needspace{7\baselineskip}\refstepcounter{ornek}\par\medskip\noindent\textbf{\color{anamavi}Örnek \theornek.}\ }{\par}
\newenvironment{cozum}{\begin{tcolorbox}[colback=cozumyesil!6,colframe=cozumyesil,boxrule=0.5pt,left=4pt,right=4pt,top=2pt,bottom=2pt,breakable]\textbf{\color{cozumyesil}Çözüm.}\ }{\end{tcolorbox}}
\newcounter{alis}
\newenvironment{alistirma}{\refstepcounter{alis}\par\medskip\needspace{9\baselineskip}\noindent\textbf{\color{anamavi}Alıştırma \thealis.}\ }{\par}
\newcommand{\kendinizoncesi}{\needspace{12\baselineskip}}
\newcommand{\ab}[1]{\overline{#1}}
\setlength{\parindent}{0pt}
\setlength{\parskip}{3pt}

\begin{document}

% ---------------- Başlık ----------------
\begin{minipage}[c]{\textwidth}
{\LARGE\bfseries\color{anamavi} MMAT105 Genel Matematik}\\[4pt]
{\large Hafta 1 Ders Notu: Temel Kavramlar, İşlem Önceliği,\\ Ardışık Sayılar ve Faktöriyel}\\[6pt]
{\small 2026--2027 Güz}
\end{minipage}

\vspace{4pt}\noindent\textcolor{anamavi}{\rule{\textwidth}{1pt}}

\textbf{Bu notun kullanımı.} Her bölümde önce kısa bir kural kutusu, ardından çözümlü örnekler yer alır. Örnekleri çözümüne bakmadan önce kendiniz deneyin. Notun sonunda önce çözümlü alıştırmalar, ardından cevap anahtarlı ``Kendiniz Çözün'' alıştırmaları yer alır.

% =====================================================
\section{Sayı Kümeleri}

\begin{kural}[Tanımlar]
\begin{itemize}[leftmargin=*,itemsep=1pt]
\item Doğal sayılar: \(\mathbb{N}=\{0,1,2,3,\dots\}\); \quad sayma sayıları: \(\mathbb{N}^{+}=\{1,2,3,\dots\}\)
\item Tam sayılar: \(\mathbb{Z}=\{\dots,-2,-1,0,1,2,\dots\}=\mathbb{Z}^{-}\cup\{0\}\cup\mathbb{Z}^{+}\), \quad \(\mathbb{N}\subset\mathbb{Z}\)
\item \(a\le b\) tam sayıları için \([a,b]\) aralığındaki tam sayı sayısı: \(b-a+1\)
\end{itemize}
\end{kural}

\begin{dikkat}
\(0\) bir doğal sayıdır; ne pozitif ne negatiftir; çift bir sayıdır.
\end{dikkat}

\begin{ornek}
\(A=\{x\in\mathbb{Z}\,:\,-3<x\le 4\}\) kümesinin kaç elemanı vardır? Bu kümenin doğal sayı olan elemanları hangileridir?
\end{ornek}
\begin{cozum}
\(-3\) dahil değil, \(4\) dahil: \(A=\{-2,-1,0,1,2,3,4\}\) \(\Rightarrow\) \(4-(-2)+1=7\) eleman.\\
Doğal sayı olanlar: \(0,1,2,3,4\) \ (\(5\) tane).
\end{cozum}

% =====================================================
\section{Rakam, Basamak ve Çözümleme}

\begin{kural}[Tanımlar ve çözümleme]
Rakamlar: \(\{0,1,\dots,9\}\). Bir rakamın basamak değeri, rakamın bulunduğu basamağın değeriyle çarpımıdır.
\[\ab{ab}=10a+b,\qquad \ab{abc}=100a+10b+c,\qquad \ab{abcd}=1000a+100b+10c+d\quad (a\neq 0)\]
Sık kullanılan sonuçlar:
\[\ab{ab}+\ab{ba}=11(a+b),\qquad \ab{ab}-\ab{ba}=9(a-b),\qquad \ab{abc}-\ab{cba}=99(a-c)\]
\end{kural}

\begin{ornek}
\(a\) ve \(b\) birbirinden farklı sıfırdan farklı rakamlardır. \(\ab{ab}+\ab{ba}\) toplamı en çok kaçtır?
\end{ornek}
\begin{cozum}
\(\ab{ab}+\ab{ba}=11(a+b)\). Farklı rakamlarla \(a+b\) en çok \(9+8=17\) olur.
\[11\cdot 17=187\qquad(\text{kontrol: } 98+89=187)\]
\end{cozum}

\begin{ornek}
\(\ab{3a}+\ab{a3}=77\) olduğuna göre \(a\) kaçtır?
\end{ornek}
\begin{cozum}
\[(30+a)+(10a+3)=33+11a=77\ \Rightarrow\ 11a=44\ \Rightarrow\ a=4\]
Kontrol: \(34+43=77\).
\end{cozum}

\begin{ornek}
İki basamaklı bir sayının sağına \(5\) yazıldığında elde edilen üç basamaklı sayı, ilk sayıdan \(419\) fazladır. İlk sayı kaçtır?
\end{ornek}
\begin{cozum}
İlk sayı \(x\) olsun. Sağına \(5\) yazmak sayıyı \(10\) ile çarpıp \(5\) eklemektir: \(10x+5\).
\[(10x+5)-x=419\ \Rightarrow\ 9x=414\ \Rightarrow\ x=46\]
Kontrol: \(465-46=419\).
\end{cozum}

% =====================================================
\section{İşlem Önceliği}

\begin{kural}[İşlem sırası]
\begin{enumerate}[leftmargin=*,itemsep=1pt]
\item Parantez içi (en içteki önce): \( (\ )\rightarrow[\ ]\rightarrow\{\ \}\)
\item Üs alma
\item Çarpma ve bölme: aynı öncelik, \textbf{soldan sağa}
\item Toplama ve çıkarma: aynı öncelik, \textbf{soldan sağa}
\end{enumerate}
Kesir çizgisi parantez gibi davranır: pay ve payda ayrı ayrı hesaplanır.
\end{kural}

\begin{dikkat}[Sık yapılan hatalar]
\(20\div 4\cdot 5=25\) (\(1\) değil); \quad \(-3^2=-9\) ama \((-3)^2=9\); \quad \(48\div 4\div 2=6\) ama \(48\div(4\div 2)=24\).
\end{dikkat}

\begin{ornek}
\(\dfrac{3^2-2\cdot(5-8)}{(-1)^3+2^2}\) işleminin sonucu kaçtır?
\end{ornek}
\begin{cozum}
Pay: \(9-2\cdot(-3)=9+6=15\). \quad Payda: \(-1+4=3\).
\[\frac{15}{3}=5\]
\end{cozum}

\begin{ornek}
\(-(-2)^3-3\cdot\big[-4+2\cdot(1-3)\big]\) işleminin sonucu kaçtır?
\end{ornek}
\begin{cozum}
\((-2)^3=-8\Rightarrow-(-8)=8\). \quad Köşeli parantez: \(-4+2\cdot(-2)=-4-4=-8\).
\[8-3\cdot(-8)=8+24=32\]
\end{cozum}

% =====================================================
\needspace{12\baselineskip}
\section{İşaret Kuralları}

\begin{kural}[İşaret]
\begin{itemize}[leftmargin=*,itemsep=1pt]
\item Çarpma/bölme: aynı işaretler \(\Rightarrow +\), farklı işaretler \(\Rightarrow -\). Tek sayıda negatif çarpan varsa sonuç negatiftir.
\item \(a>0\) için \((-a)^{2k}>0\), \((-a)^{2k+1}<0\); \quad \((-1)^{\text{çift}}=1\), \((-1)^{\text{tek}}=-1\).
\item \(a-(-b)=a+b\).
\end{itemize}
\end{kural}

\begin{ornek}
\(a<b<0<c\) olduğuna göre \(\dfrac{a\cdot c}{b}\) ve \((a-b)\cdot c\) ifadelerinin işaretleri nedir?
\end{ornek}
\begin{cozum}
\(a\cdot c\): \((-)(+)=-\); \ bunu \(b<0\) ile bölmek: \((-)/(-)=+\) \(\Rightarrow\) \(\dfrac{a\cdot c}{b}\) \textbf{pozitif}.\\
\(a<b\Rightarrow a-b<0\); \ \((-)\cdot(+)=-\) \(\Rightarrow\) \((a-b)\cdot c\) \textbf{negatif}.
\end{cozum}

\begin{ornek}
\(n\) bir pozitif tam sayıdır. \((-1)^{n}+(-1)^{n+1}+(-1)^{2n}\) ifadesinin değeri kaçtır?
\end{ornek}
\begin{cozum}
\(n\) ile \(n+1\) ardışık olduğundan biri tek, diğeri çifttir: \((-1)^{n}+(-1)^{n+1}=1+(-1)=0\) \ (sıra fark etmez).\\
\(2n\) her zaman çifttir: \((-1)^{2n}=1\). \quad Sonuç: \(0+1=1\).
\end{cozum}

% =====================================================
\section{Tek ve Çift Sayılar}

\begin{kural}[Tek -- çift]
Çift: \(2k\), \ tek: \(2k+1\) \((k\in\mathbb{Z})\).
\begin{center}
\begin{tabular}{c|cc}
\(\pm\) & T & Ç\\\hline T & Ç & T\\ Ç & T & Ç
\end{tabular}
\qquad\qquad
\begin{tabular}{c|cc}
\(\cdot\) & T & Ç\\\hline T & T & Ç\\ Ç & Ç & Ç
\end{tabular}
\end{center}
Çarpanlardan en az biri çiftse çarpım çifttir. Çarpım tekse bütün çarpanlar tektir. \(n\) tek ise \(n^{k}\) de tektir.
\end{kural}

\begin{ornek}
\(a\) ve \(b\) tam sayılardır. \(a\) çift ve \(a^2+3b\) tek ise \(b\) tek midir, çift midir? \(a\cdot b+b\) tek midir, çift midir?
\end{ornek}
\begin{cozum}
\(a\) çift \(\Rightarrow a^2\) çift. \(\text{Ç}+3b=\text{T}\Rightarrow 3b\) tek \(\Rightarrow b\) \textbf{tek}.\\
\(a\cdot b\): Ç\(\cdot\)T \(=\) Ç; \ \(\text{Ç}+b=\text{Ç}+\text{T}=\) \textbf{tek}.
\end{cozum}

% =====================================================
\section{Ardışık Sayılar}

\begin{kural}[Terim sayısı ve toplam]
Artış miktarı sabit olan sayılar için:
\[\text{Terim sayısı}=\frac{\text{Son terim}-\text{İlk terim}}{\text{Artış}}+1,\qquad \text{Toplam}=\text{Terim sayısı}\cdot\frac{\text{İlk terim}+\text{Son terim}}{2}\]
Özel toplamlar:
\[1+2+\dots+n=\frac{n(n+1)}{2},\qquad 2+4+\dots+2n=n(n+1),\qquad 1+3+\dots+(2n-1)=n^2\]
Tek sayıda ardışık terimin toplamı \(=\) terim sayısı \(\times\) ortanca terim.
\end{kural}

\begin{ornek}
\(7+12+17+\dots+152\) toplamında kaç terim vardır? Toplam kaçtır?
\end{ornek}
\begin{cozum}
Artış \(5\): \quad \(n=\dfrac{152-7}{5}+1=29+1=30\).
\[\text{Toplam}=30\cdot\frac{7+152}{2}=15\cdot 159=2385\]
\end{cozum}

\begin{ornek}
\(1\)'den \(40\)'a kadar olan çift sayıların toplamı ile tek sayıların toplamı arasındaki fark kaçtır?
\end{ornek}
\begin{cozum}
Çiftler: \(2+4+\dots+40\), \(n=20\): \(20\cdot 21=420\).\\
Tekler: \(1+3+\dots+39\), \(n=20\): \(20^2=400\).
\[420-400=20\]
Kısa yol: \((2-1)+(4-3)+\dots+(40-39)\) \(\Rightarrow\) \(20\) tane \(1\).
\end{cozum}

\begin{ornek}
Ardışık \(4\) tek sayının toplamı \(72\) ise en küçüğü kaçtır?
\end{ornek}
\begin{cozum}
Sayılar \(x,\ x+2,\ x+4,\ x+6\):
\[4x+12=72\ \Rightarrow\ x=15\]
Sayılar: \(15,17,19,21\). (Kontrol: toplam \(72\).)
\end{cozum}

\begin{ornek}
Üç basamaklı doğal sayılardan kaç tanesi \(7\) ile tam bölünür?
\end{ornek}
\begin{cozum}
En küçük: \(7\cdot 15=105\); \ en büyük: \(7\cdot 142=994\).
\[n=\frac{994-105}{7}+1=127+1=128\qquad(\text{ya da } 142-15+1=128)\]
\end{cozum}

% =====================================================
\section{Faktöriyel}

\begin{kural}[Tanım ve özellikler]
\[n!=1\cdot 2\cdot 3\cdots n\quad(n\in\mathbb{N}^{+}),\qquad 0!=1,\qquad n!=n\cdot(n-1)!\]
\begin{center}
\begin{tabular}{cccccccc}
\toprule
\(1!\) & \(2!\) & \(3!\) & \(4!\) & \(5!\) & \(6!\) & \(7!\) & \(8!\)\\\midrule
\(1\) & \(2\) & \(6\) & \(24\) & \(120\) & \(720\) & \(5040\) & \(40320\)\\\bottomrule
\end{tabular}
\end{center}
\textbf{Sadeleştirme:} Küçük faktöriyel ortak çarpan olarak dışarı alınır: \(7!+6!=6!\,(7+1)=8\cdot 6!\).\\
\textbf{Sondaki sıfır sayısı:} \(n!\) sayısının sonundaki sıfır sayısı \(\left\lfloor\frac{n}{5}\right\rfloor+\left\lfloor\frac{n}{25}\right\rfloor+\left\lfloor\frac{n}{125}\right\rfloor+\cdots\) \ (\(\lfloor x\rfloor\): \(x\)'in tam kısmı)
\end{kural}

\begin{dikkat}
\(5!\) ve sonraki bütün faktöriyellerin birler basamağı \(0\)'dır. \quad \((a+b)!\neq a!+b!\), \quad \((2n)!\neq 2\cdot n!\).
\end{dikkat}

\begin{ornek}
\(\dfrac{9!}{7!\cdot 3!}\) ve \(\dfrac{6!+7!+8!}{6!}\) işlemlerinin sonuçları nedir?
\end{ornek}
\begin{cozum}
\[\frac{9!}{7!\cdot 3!}=\frac{9\cdot 8\cdot 7!}{7!\cdot 6}=\frac{72}{6}=12\]
\[\frac{6!+7!+8!}{6!}=\frac{6!\,(1+7+7\cdot 8)}{6!}=1+7+56=64\]
\end{cozum}

\begin{ornek}
\(n!+(n-1)!=7\cdot(n-1)!\) olduğuna göre \(n\) kaçtır?
\end{ornek}
\begin{cozum}
\(n!=n\cdot(n-1)!\) yazılırsa:
\[n\cdot(n-1)!+(n-1)!=(n-1)!\,(n+1)=7\cdot(n-1)!\ \Rightarrow\ n+1=7\ \Rightarrow\ n=6\]
\end{cozum}

\begin{ornek}
\(50!\) sayısının sonunda kaç sıfır vardır?
\end{ornek}
\begin{cozum}
Her sıfır bir \(2\cdot 5\) çiftinden gelir; \(2\) çarpanı \(5\)'ten fazla olduğundan \(5\) çarpanlarını saymak yeterlidir.
\[\left\lfloor\frac{50}{5}\right\rfloor+\left\lfloor\frac{50}{25}\right\rfloor=10+2=12\]
(\(25\) ve \(50\) sayıları ikişer tane \(5\) çarpanı içerir; ikinci terim bunları sayar.)
\end{cozum}

\begin{ornek}
\(1!+2!+3!+\dots+10!\) toplamının birler basamağındaki rakam kaçtır?
\end{ornek}
\begin{cozum}
\(5!\)'den itibaren her terimin birler basamağı \(0\)'dır. Sadece ilk dört terim önemlidir:
\[1+2+6+24=33\ \Rightarrow\ \text{birler basamağı } 3\]
\end{cozum}

% =====================================================
\section{Özet: Formül Kartı}
\begin{center}
\renewcommand{\arraystretch}{1.35}
\begin{tabular}{@{}ll@{}}
\toprule
Konu & Formül / Kural\\\midrule
Aralıktaki tam sayı sayısı & \([a,b]\): \(b-a+1\)\\
Çözümleme & \(\ab{ab}\pm\ab{ba}\): \(11(a+b)\), \ \(9(a-b)\); \ \(\ab{abc}-\ab{cba}=99(a-c)\)\\
İşlem sırası & parantez \(\to\) üs \(\to\) \(\times,\div\) \(\to\) \(+,-\) (soldan sağa)\\
İşaret & tek sayıda negatif çarpan \(\Rightarrow\) negatif; \ \(-a^2\neq(-a)^2\)\\
Tek -- çift & bir çarpan çift \(\Rightarrow\) çarpım çift; \ T\(\pm\)T\(=\)Ç\\
Terim sayısı & \(\dfrac{\text{son}-\text{ilk}}{\text{artış}}+1\)\\
Ardışık toplam & terim sayısı \(\cdot\dfrac{\text{ilk}+\text{son}}{2}\); \ \(\dfrac{n(n+1)}{2}\), \(n(n+1)\), \(n^2\)\\
Faktöriyel & \(n!=n\,(n-1)!\), \ \(0!=1\); \ sondaki sıfır: \(\lfloor n/5\rfloor+\lfloor n/25\rfloor+\cdots\)\\\bottomrule
\end{tabular}
\end{center}

% =====================================================
\needspace{20\baselineskip}
% Hafta 1 -- alıştırma gövdesi (ders notu ve ayrı alıştırma kağıdı tarafından \input edilir)
\section{Çözümlü Alıştırmalar}
Her soruyu önce kendiniz deneyin, sonra çözümü okuyun.

\begin{alistirma}
Rakamları birbirinden farklı üç basamaklı en büyük çift doğal sayı ile rakamları birbirinden farklı üç basamaklı en küçük tek doğal sayının farkı kaçtır?
\end{alistirma}
\begin{cozum}
\textbf{En büyük çift sayı:} Yüzler basamağına en büyük rakam \(9\), onlara \(8\) yazılır. Birler basamağı çift olmalı ve \(9,8\)'den farklı olmalı: en büyük seçenek \(6\). \(\Rightarrow 986\).\\
\textbf{En küçük tek sayı:} Yüzler basamağı \(0\) olamaz \(\Rightarrow 1\); onlar \(0\). Birler basamağı tek olmalı ve \(1\)'den farklı olmalı: en küçük seçenek \(3\). \(\Rightarrow 103\).
\[986-103=883\]
\end{cozum}

\begin{alistirma}
\(\ab{abc}\) ve \(\ab{cba}\) üç basamaklı sayılardır. \(\ab{abc}-\ab{cba}=594\) ve \(a+c=12\) olduğuna göre \(a\cdot c\) kaçtır?
\end{alistirma}
\begin{cozum}
Çözümleme: \((100a+10b+c)-(100c+10b+a)=99a-99c=99(a-c)\).
\[99(a-c)=594\ \Rightarrow\ a-c=6\]
\(a-c=6\) ve \(a+c=12\) taraf tarafa toplanırsa \(2a=18\Rightarrow a=9\), \ \(c=3\).
\[a\cdot c=9\cdot 3=27\qquad(\text{kontrol, } b=5 \text{ için: } 953-359=594)\]
\end{cozum}

\begin{alistirma}
\(\dfrac{(-2)^3+4^2\div 2}{-3^2+10}\) işleminin sonucu kaçtır?
\end{alistirma}
\begin{cozum}
Kesir çizgisi parantez gibidir; pay ve payda ayrı hesaplanır.\\
Pay: \((-2)^3=-8\) (tek kuvvet, işaret korunur); \ \(4^2\div 2=16\div 2=8\); \ \(-8+8=0\).\\
Payda: \(-3^2=-9\) (üs yalnız \(3\)'e ait); \ \(-9+10=1\).
\[\frac{0}{1}=0\]
\end{cozum}

\begin{alistirma}
\(x,y,z\) tam sayılardır. \(x\cdot y\) tek ve \(x+z\) çift ise \(z\) tek midir, çift midir? \(x+y+z\) tek midir, çift midir?
\end{alistirma}
\begin{cozum}
Çarpım tek ise bütün çarpanlar tektir: \(x\) tek, \(y\) tek.\\
\(x+z\) çift ve \(x\) tek \(\Rightarrow\) \(z\) de tek olmalı (T\(+\)T\(=\)Ç). \quad \(z\) \textbf{tek}.\\
\(x+y+z=\) T\(+\)T\(+\)T \(=\) Ç\(+\)T \(=\) \textbf{tek}. \quad (Kontrol: \(1+3+5=9\).)
\end{cozum}

\begin{alistirma}
\(13+17+21+\dots+97\) toplamında kaç terim vardır? Toplam kaçtır?
\end{alistirma}
\begin{cozum}
Artış \(17-13=4\).
\[\text{Terim sayısı}=\frac{97-13}{4}+1=\frac{84}{4}+1=21+1=22\]
\[\text{Toplam}=22\cdot\frac{13+97}{2}=22\cdot 55=1210\]
\end{cozum}

\begin{alistirma}
Ardışık \(5\) tam sayının toplamı \(-15\) ise bu sayıların en büyüğü kaçtır?
\end{alistirma}
\begin{cozum}
Terim sayısı tek olduğundan toplam \(=5\cdot\) ortanca terim:
\[\text{Ortanca}=\frac{-15}{5}=-3\]
Sayılar: \(-5,\ -4,\ -3,\ -2,\ -1\). \quad En büyüğü \(-1\).\\
(Dikkat: Negatif sayılarda en büyük, sıfıra en yakın olandır.)
\end{cozum}

\begin{alistirma}
\(n\) bir doğal sayıdır. \(\dfrac{(n+2)!}{n!}=56\) olduğuna göre \(n\) kaçtır?
\end{alistirma}
\begin{cozum}
\((n+2)!=(n+2)(n+1)\cdot n!\) yazılırsa:
\[\frac{(n+2)(n+1)\cdot n!}{n!}=(n+2)(n+1)=56\]
Ardışık iki doğal sayının çarpımı \(56=7\cdot 8\) \(\Rightarrow\) \(n+1=7\), \(n=6\).\\
(Kontrol: \(8!/6!=8\cdot 7=56\).)
\end{cozum}

\begin{alistirma}
\(2!+4!+6!+\dots+20!\) toplamının birler basamağındaki rakam kaçtır? Bu toplam tek midir, çift midir?
\end{alistirma}
\begin{cozum}
\(n\ge 5\) için \(n!\) içinde \(2\cdot 5=10\) çarpanı bulunur; birler basamağı \(0\)'dır. Bu yüzden \(6!,8!,\dots,20!\) birler basamağını etkilemez.\\
Kalan terimler: \(2!+4!=2+24=26\) \(\Rightarrow\) birler basamağı \(6\).\\
Birler basamağı çift olduğundan toplam \textbf{çifttir} (zaten her terim çifttir).
\end{cozum}

\kendinizoncesi
\section{Alıştırmalar (Kendiniz Çözün)}

\begin{enumerate}[label=\textbf{\arabic*.},leftmargin=*,itemsep=6pt]
\item \(B=\{x\in\mathbb{N}\,:\,x<5\}\) ve \(C=\{x\in\mathbb{Z}\,:\,-2\le x<3\}\) kümelerinin eleman sayıları kaçtır?
\item \(-7\) ile \(6\) arasında (uçlar hariç) kaç tam sayı vardır?
\item \(\ab{ab}\) ve \(\ab{ba}\) iki basamaklı sayılardır. \(\ab{ab}+\ab{ba}=143\) ve \(a-b=5\) olduğuna göre \(\ab{ab}\) sayısı kaçtır?
\item İki basamaklı \(\ab{ab}\) sayısı, rakamları toplamının \(7\) katına eşittir. Bu koşulu sağlayan kaç sayı vardır?
\item \(36\div 6\cdot 2-3\cdot(2-5)^2\) işleminin sonucu kaçtır?
\item \(-2^4+(-2)^3-(-1)^{99}\) işleminin sonucu kaçtır?
\item \(a\cdot b>0\), \ \(b\cdot c<0\) ve \(a<0\) olduğuna göre \(b\) ve \(c\)'nin işaretleri nedir? \(a\cdot c\) pozitif midir, negatif midir?
\item \(a\) ve \(b\) tam sayılardır. \(a\) tek ve \(a^2\cdot b+a\) çift ise \(b\) tek midir, çift midir?
\item \(1+3+5+\dots+59\) ve \(2+4+6+\dots+60\) toplamları kaçtır?
\item Ardışık \(3\) tek sayının toplamı \(117\) ise en büyüğü kaçtır?
\item İki basamaklı doğal sayılardan kaç tanesi \(6\)'nın katıdır?
\item \(11+12+13+\dots+40\) toplamı kaçtır?
\item \(\dfrac{10!}{8!}-\dfrac{7!}{5!}\) işleminin sonucu kaçtır?
\item \(\dfrac{7!+8!}{9\cdot 6!}\) işleminin sonucu kaçtır?
\item \(40!\) sayısının sonunda kaç sıfır vardır?
\end{enumerate}

\subsection*{Cevap Anahtarı}
{\small
\begin{multicols}{2}
\begin{enumerate}[label=\textbf{\arabic*.},leftmargin=*,itemsep=1pt]
\item \(B=\{0,1,2,3,4\}\): \(5\); \ \(C=\{-2,\dots,2\}\): \(5\).
\item \(-6,\dots,5\): \(5-(-6)+1=12\).
\item \(a+b=13\), \(a-b=5\Rightarrow a=9,\ b=4\); \ \(94\).
\item \(10a+b=7a+7b\Rightarrow a=2b\): \(21,42,63,84\); \ \(4\) tane.
\item \(6\cdot 2-3\cdot 9=12-27=-15\).
\item \(-16+(-8)-(-1)=-23\).
\item \(b<0\), \(c>0\); \ \(a\cdot c\) negatif.
\item \(a^2\) tek; \(a^2b\) tek olmalı \(\Rightarrow\) \(b\) tek.
\item \(30^2=900\); \ \(30\cdot 31=930\).
\item Ortanca \(39\): \(37,39,41\); \ en büyük \(41\).
\item \(12,18,\dots,96\): \(\frac{96-12}{6}+1=15\).
\item \(\frac{40\cdot 41}{2}-\frac{10\cdot 11}{2}=820-55=765\).
\item \(10\cdot 9-7\cdot 6=90-42=48\).
\item \(\frac{6!\,(7+56)}{9\cdot 6!}=\frac{63}{9}=7\).
\item \(\lfloor 40/5\rfloor+\lfloor 40/25\rfloor=8+1=9\).
\end{enumerate}
\end{multicols}}


\vspace{6pt}
\noindent\textcolor{anamavi}{\rule{\textwidth}{0.5pt}}\\
{\small Gelecek hafta: Asal sayılar, asal çarpanlara ayırma, tam bölenler, bölme.}

\end{document}
